\chapter{Learning paths}\label{sec:root:learning-paths:1}

Not every reader needs every chapter in order. This page maps out how the chapters actually depend on one another, and suggests a handful of named paths through the book for common starting points. All of them converge by Chapter 6 --- the dependency structure only really branches in Part I.

\section{Chapter dependency graph}\label{sec:root:learning-paths:2}

\begin{center}
% Chapter dependency graph: a single chain, Chapter 0 through 13, with the
% two checkpoint projects as waypoints (diamond nodes) after Chapters 5
% and 11. Source: learning-paths.md.
\begin{tikzcd}[column sep=small, row sep=small]
\text{0. Setup} \arrow[d] \\
\text{1. Basics} \arrow[d] \\
\text{2. Functions \& structures} \arrow[d] \\
\text{3. Propositions \& proofs} \arrow[d] \\
\text{4. Tactics} \arrow[d] \\
\text{5. Rigor check} \arrow[d] \\
\langle\text{Checkpoint: Monoid}\rangle \arrow[d] \\
\text{6. Groups} \arrow[d] \\
\text{7. Group theorems} \arrow[d] \\
\text{8. Rings} \arrow[d] \\
\text{9. Ring theorems} \arrow[d] \\
\text{10. Modules} \arrow[d] \\
\text{11. Path algebras} \arrow[d] \\
\langle\text{Checkpoint: Path.length}\rangle \arrow[d] \\
\text{12. Working efficiently} \arrow[d] \\
\text{13. Next steps}
\end{tikzcd}

\end{center}

Solid arrows are hard prerequisites (each chapter's Lean code and proofs genuinely build on the one before it). The two checkpoint projects are optional but recommended waypoints, not prerequisites --- nothing later in the book requires having done them. Dashed arrows are the two named paths below that actually skip material outright, rather than just reading it faster; the other two named paths change \emph{how} a chapter is read, not which chapters are read, so they have no edge of their own.

Chapter 5 is the one true fork: it exists to answer rigor questions (\texttt{structure} vs \texttt{class}, universes, definitional vs propositional equality) that a careful reader will already be asking by Chapter 4, but a reader willing to take Lean's guarantees on faith for now can skip it on a first pass and come back once something in Chapter 6 onward prompts the question directly (each chapter's ``Read more'' boxes still point back to the relevant §).

\section{Named paths}\label{sec:root:learning-paths:3}

\textbf{Full path (recommended for a first read).} Chapters 0--13 in order, checkpoint projects included. This is the path the book is written to support directly --- every forward reference assumes it.

\textbf{``I already know Lean, teach me the algebra.''} Skim Chapter 0 (just confirm your toolchain matches \texttt{v4.31.0}), read Chapter 1 §1--§3 for this book's specific \texttt{Fin}/\texttt{Vec} examples, skip §4--§5 and Chapter 5 entirely unless something later sends you back (the \hyperref[sec:01-basics:04-terminology:1]{Chapter 1 §4 glossary} and \hyperref[sec:root:tactic-and-library-reference:1]{tactic and library reference} work as pure lookup tables if a term is unfamiliar), then read Chapters 2--4 quickly for this book's own conventions before Chapters 6--13 in full.

\textbf{``I already know abstract algebra, teach me Lean.''} Read Chapters 0--5 in full --- this is the actual Lean-specific content, and none of it assumes algebra beyond what you already have. In Chapters 6--11, skim the mathematical statements (you already know why the theorems are true) and concentrate on \emph{how} each is expressed and proved in Lean, plus the ``Mathlib equivalent'' boxes; read Chapters 12--13 in full.

\textbf{``I want the formal foundations before anything else.''} Read Chapter 1 in full, including §4--§5's calculus-of-constructions material, then Chapter 3 §1--§2 for the logic recap, then Chapter 5 in full including §3's typing rules --- all before touching tactics or groups. This front-loads everything Chapters 6 onward take for granted, at the cost of deferring concrete payoff the longest.

\textbf{``I want to see Lean do real mathematics as fast as possible.''} Read Chapter 0 §1 (``Why Lean?''), then jump straight to Chapters 6--7 (\texttt{Group}), referring back to Chapters 1--5 only when a specific term or tactic is unfamiliar (the \hyperref[sec:01-basics:04-terminology:1]{Chapter 1 §4 glossary} and

built exactly for this kind of lookup). Continue to Chapters 8--11, then loop back and fill in whichever of Chapters 1--5 you skipped once you have enough concrete motivation to want the ``why.''

Whichever path is chosen, the two checkpoint projects (after Chapter 5, after Chapter 11) are a good self-test of whether to continue forward or double back first.
